% ATTEMPT_STATUS: unresolved
% RESEARCH_GENERATED: 2026-10-01; GPT-6 Astra Ultra; individual mathematical attempt
% TOP_PROBLEM: {"id": 168, "title": "Cap Sets in F_7^n", "category_id": 2, "category": {"id": 2, "name": "combinatorics", "display_name": "Combinatorics", "description": "Counting problems, graph theory, discrete structures.", "slug": "combinatorics", "order_index": 2, "created_at": "2026-07-31T15:26:25.671Z"}, "status": "open"}
% FIRST_POSTED: 2026-10-01
% Imported from: unsolvedmath_additional_500.tex; UnsolvedMath ID 168
% !TeX program = lualatex
% Compile twice: lualatex 168.tex
% Self-contained: no external bibliography, images, or \input files.
% Numeric section labels and visible numbers are original UnsolvedMath IDs.
\documentclass[11pt,a4paper]{article}
\usepackage[margin=24mm,headheight=14pt]{geometry}
\usepackage{fontspec}
\setmainfont{Latin Modern Roman}
\setsansfont{Latin Modern Sans}
\setmonofont{Latin Modern Mono}
\usepackage{amsmath,amssymb,mathtools}
\usepackage{unicode-math}
\setmathfont{Latin Modern Math}
\usepackage{microtype}
\usepackage{enumitem,xurl,needspace}
\usepackage[dvipsnames]{xcolor}
\usepackage{hyperref}
\hypersetup{unicode=true,colorlinks=true,linkcolor=MidnightBlue,urlcolor=MidnightBlue,
 pdftitle={UnsolvedMath Problem 168},pdfauthor={Source-annotated compilation},bookmarksnumbered=true}
\usepackage{bookmark,tocloft,titlesec,fancyhdr}
\setlength{\cftsecnumwidth}{4.2em}
\cftsetpnumwidth{2.5em}
\cftsetrmarg{4em}
\titleformat{\section}{\Large\bfseries\raggedright}{\thesection}{0.7em}{}
\pagestyle{fancy}\fancyhf{}
\fancyhead[L]{\small\sffamily UnsolvedMath research attempt}
\fancyhead[R]{\small Original catalogue IDs}
\fancyfoot[C]{\thepage}
\setlength{\parindent}{0pt}
\setlength{\parskip}{5pt plus 1pt}
\setlength{\emergencystretch}{3em}
\setcounter{tocdepth}{1}\setcounter{secnumdepth}{1}
\setlist[itemize]{leftmargin=3em,itemsep=4pt,topsep=4pt,parsep=0pt}
\allowdisplaybreaks
\newcommand{\N}{\mathbb{N}}
\newcommand{\Z}{\mathbb{Z}}
\newcommand{\Q}{\mathbb{Q}}
\newcommand{\R}{\mathbb{R}}
\newcommand{\C}{\mathbb{C}}
\newcommand{\F}{\mathbb{F}}
\newcommand{\A}{\mathbb{A}}
\newcommand{\PP}{\mathbb{P}}
\newcommand{\E}{\mathbb{E}}
\newcommand{\eps}{\varepsilon}
\newcommand{\dd}{\,\mathrm d}
\newcommand{\sref}[2]{\hyperlink{source:#1:#2}{[S#2]}}
\newif\ifshowcatalogue
\showcataloguetrue
\newcommand{\cataloguescope}[1]{\ifshowcatalogue
 \paragraph{Statement recorded in the supplied source.}
 #1\fi}
\begin{document}
% UnsolvedMath ID 168; source catalogue code GREEN-080

\setcounter{section}{167}

\section{Avoiding Small-Coordinate Differences over F\_7}\label{168}

{\small\sffamily UnsolvedMath ID 168 · Combinatorics}

\subsection{Definitions and mathematical statement}

\paragraph{Shared notation.}Unless a section says otherwise, $\N=\{1,2,\ldots\}$,
$\Z$ is the ring of integers, $\Q,\R,\C$ are the rational, real, and complex
fields, $[N]=\{1,\ldots,\lfloor N\rfloor\}$, and $\log$ is the natural
logarithm. For real functions of a parameter tending to infinity,
$f\ll g$ and $f=O(g)$ mean $|f|\leq Cg$ eventually for some fixed $C$;
$f\gg g$ reverses this comparison; $f=o(g)$ means $f/g\to0$;
$f\sim g$ means $f/g\to1$; and $f\asymp g$ means both order bounds.
A subscript on $O$ or $\ll$ specifies allowed constant dependence.
The expression $n^{o(1)}$ in an upper bound means growth smaller than
$n^\epsilon$ for each fixed $\epsilon>0$. Local definitions govern any
exception, finite-field convention, or use of the same symbol for another
object. An asymptotic claim is not automatically an effective algorithm.



Represent $-1,0,1$ by their residues in $\F_7$. Set
\[M(n)=\max\{|A|:A\subseteq\F_7^n,\ (A-A)\cap\{-1,0,1\}^n=\{0\}\}.\]
Thus no two distinct chosen vectors may differ only by coordinates in $\{-1,0,1\}$. This is the precise difference-avoidance property, not the usual three-term-progression definition of a cap set. \sref{168}{1} \sref{168}{2}

\paragraph{Statement recorded in the supplied source.}
What is the largest subset $A \subset \mathbb{F}_7^n$ for which $A - A$ intersects $\{-1, 0, 1\}^n$ only at 0? \sref{168}{1}

\paragraph{Residual task reported by the archive.}

The embedded review (2026-08-17) records: Close the gap between the 367\textasciicircum{}(1/5) construction and the Lovász theta upper bound, or otherwise determine Theta(C\_7). \sref{168}{1}

\subsection{Short English statement}

How many vectors over the seven-element field can be chosen so that no two have all coordinate differences equal to zero or plus or minus one?

\subsection{Research attempt}

\paragraph{Provenance and scope.}
This individual attempt was generated by GPT-6 Astra Ultra on 1 October 2026. The method was to derive explicit partial results and test the first proposed extension adversarially. The original statement and sources below are retained. No claim of mathematical novelty or complete solution is made.

\paragraph{Formal target and primary-source comparison.}
Let $M(n)$ be the maximum cardinality of $A\subset\F_7^n$ such that no nonzero difference lies in $\{-1,0,1\}^n$. This is a strong-product cycle independence problem, not a cap-set condition. Polak and Schrijver's inspected paper supplies a much stronger five-coordinate construction of size $367$ than the small block below. We make no record claim. The aim is an explicitly checked two-coordinate construction, a packing upper bound, and a careful separation between finite-dimensional maxima and their asymptotic exponential rate.

\paragraph{A disjoint-box upper bound.}
For each $a\in A$, consider the translate $a+\{0,1\}^n$. These boxes are pairwise disjoint. Indeed, if $a+u=b+v$, then $a-b=v-u\in\{-1,0,1\}^n$. The avoidance property gives $a=b$, and then $u=v$. Each box contains $2^n$ points, while the ambient space has $7^n$. Hence
\[
 M(n)\leq(7/2)^n.
\]
There is no Euclidean-volume approximation in this proof: it is exact counting in the finite group. The box choice is important because its difference set is precisely the forbidden coordinate cube.

\paragraph{A verified ten-point block.}
In two coordinates, take the following row supports, with coordinates read modulo seven:
\[
 \begin{array}{c|c}
 x&\{y:(x,y)\in B\}\\\hline
 0&\{1\}\\1&\{4,6\}\\2&\{2\}\\3&\{0\}\\
 4&\{3,5\}\\5&\{1\}\\6&\{4,6\}
 \end{array}
\]
The total is ten points. Within each row containing two points, the vertical difference is two modulo seven, so it is allowed. Between nonadjacent distinct rows, the horizontal difference already lies outside $\{-1,0,1\}$, so no vertical check is needed. Between adjacent rows, including rows six and zero, the listed vertical supports have every difference outside $\{-1,0,1\}$. Checking the seven neighbouring row pairs verifies the entire condition. The companion script independently tests all forty-five point pairs by exact modular arithmetic.

This construction was found by a fixed-seed finite greedy search and then verified; the search is not an exhaustive proof that ten is the two-dimensional maximum. The certificate itself, the displayed set and its pairwise check, is independent of how it was found.

\paragraph{Tensorisation and an actual limit.}
If $A\subset\F_7^r$ and $B\subset\F_7^s$ satisfy the condition, so does $A\times B$. A forbidden nonzero difference in the product would project to forbidden differences in both blocks; each projection must vanish, contradicting nonzero total difference. Therefore
\[
 M(r+s)\geq M(r)M(s).
\]
Using the ten-point block gives $M(2m)\geq10^m$. In one coordinate, $\{0,2,4\}$ is admissible, so $M(2m+1)\geq3\cdot10^m$. The resulting exponential base is $\sqrt{10}$, an elementary improvement over the coordinatewise base-three construction but weaker than the cited base $367^{1/5}$.

For clarity, the limit $\lim_{n\to\infty}M(n)^{1/n}$ exists. Set $a_n=\log M(n)$. Supermultiplicativity makes $a_n$ superadditive, and the packing bound gives $a_n/n\leq\log(7/2)$. For any fixed block length $r$, write $n=kr+s$ with $0\leq s<r$. Then $a_n\geq ka_r+a_s\geq ka_r$, so $\liminf a_n/n\geq a_r/r$. Taking the supremum over $r$ matches the obvious upper bound on the limsup. This proves the limit directly without confusing a supremum of block rates with an unproved convergence claim.

\paragraph{Why a finite block is not an upper-bound theorem.}
Even an exact computation of $M(r)$ at one or several small dimensions supplies tensor lower bounds. It does not imply $M(n)\leq M(r)^{n/r}$: the product inequality goes in the opposite direction. Larger blocks may encode interactions unavailable to any smaller tensor construction, as the five-coordinate result illustrates. The disjoint-box bound also leaves substantial unused structure and is not asserted to match the sharper known analytic bounds.

\paragraph{Verification and final gap.}
All residues are reduced modulo seven, including the wraparound adjacency between zero and six. The boxes have genuinely distinct points because the field has characteristic seven. The work proves a finite block certificate, a universal packing bound, and existence of the asymptotic rate. It does not determine that rate or the exact sequence $M(n)$. The unresolved step is a matching upper method or an improved family of blocks that closes the established literature gap.

\subsection{Sources}

{\small

\begin{itemize}

\item[\hypertarget{source:168:1}{S1}] ULAM.AI, \emph{UnsolvedMath}, supplied \texttt{problems.json}, record 168 (GREEN-080), statement and background. The embedded literature-review date is 2026-08-17; that is the archive’s review date, not a fresh certification. Dataset landing page: \url{https://huggingface.co/datasets/ulamai/UnsolvedMath}.

\item[\hypertarget{source:168:2}{S2}] S. Polak, A. Schrijver, New lower bound on the Shannon capacity of C7 from circular graphs, Information Processing Letters 143 (2019), 37-40; arXiv:1808.07438. \url{https://arxiv.org/abs/1808.07438}. Bibliographic information imported from the supplied record.

\item[\hypertarget{source:168:3}{S3}] B. Green, 100 Open Problems, Problem 38, version. \url{https://people.maths.ox.ac.uk/greenbj/papers/open-problems.pdf}. The author-hosted document was inspected on 27 September 2026; the archive’s local code is not a problem-number locator in this paper.

\item[E1] S. Polak and A. Schrijver, \emph{New lower bound on the Shannon capacity of $C_7$ from circular graphs}. \url{https://arxiv.org/abs/1808.07438}. Inspected 1 October 2026.

\end{itemize}

}

\label{end:168}
\end{document}
